\documentclass[11pt,a4paper]{article}
\usepackage[utf8]{inputenc}
\usepackage[T1]{fontenc}
\usepackage[spanish]{babel}
\usepackage{hyperref}
\usepackage{geometry}
\geometry{margin=2.5cm}
\usepackage{parskip}

\title{069 si no se llenan los explicitos en un juguete no hay limites establecidos para lo}
\author{Alejandro Valenzuela Rivera\\\small Universidad Autónoma de Nuevo León}
\date{\today}
\begin{document}
\maketitle

\begin{quote}
\textbf{Sesgo declarado.} Este trabajo se sostiene desde el sesgo declarado de su autor, y no finge neutralidad. Lo que afirma se puede tumbar, y ese es el punto: el falsificador está escrito, no escondido.
\end{quote}

\section*{Resumen}
Un campo explicito vacio no se queda vacio: lo llena quien lo use, con su criterio. Y asi un juguete que parecia inofensivo se vuelve un sistema sin limites, porque el que lo llena pone los suyos y nadie los declara.

\section{Por qué se sostiene}
Es la razon entera de que haya explicitos: no son burocracia, son los limites que el autor pone antes de que otro los ponga por el. Sin ellos, lo que se crea con el juguete tiene los limites del que lo llene, y esos limites no se ven.

\section{De dónde viene}
La autoridad, en la conversacion del 2026-10-25.

\section{El caso}
El juguete: una cosa hecha para que otro la use sin saber como funciona.

\section{Método}
Un juguete con sus limites escritos antes de usarse.

\section{Límites, y qué lo tumba}
Un juguete sin limites explicitos, con el criterio del que lo llena.

\section{Lo que queda abierto}
\textquestiondown{}Si el juguete no trae sus limites escritos, de quien son los limites que tendra?

\section*{Cómo se leyó esta tesis}
Esto se escribió en el repositorio man-made-empathy, del que este candado es una pieza: candados/06 del caso/069 si no se llenan los explicitos en un juguete no hay limites establecidos para lo.md. La máquina que lo compuso lo hizo desde este mismo archivo, y no agregó nada que no estuviera escrito. Si algo de esto es falso, se dice en el mismo sitio, y se corrige con fecha.

\end{document}
